\section{Step 30 Review the observable invariants}\label{maze-s30}

The algorithm, its picture, and its tests share explicit meanings for each field. Those meanings remain useful when the implementation changes.

\subsection{Write or inspect}

Explain queue uniqueness, the parent-distance relation, and the final path count without reading the reference source. Then change the neighbor order and rerun the checks.

\subsection{Expected observation}

At every phase-one checkpoint, the queue history contains each discovered cell once. Each discovered non-start cell has an adjacent predecessor one level nearer the start. A successful final mask is exactly one shortest chain.

\subsection{Check your understanding}

What remains valid when neighbor order changes?

\begin{itemize}
\item \textbf{A} The exact parent of every discovered cell.
\item \textbf{B} The exact whole-program instruction count.
\item \textbf{C} Queue uniqueness and shortest-path validity.
\end{itemize}
